\documentclass[12pt]{article}

% =======================================================
% 1. ENGINE + TYPOGRAPHY
%
% Font presets are controlled by the 'fontset' YAML variable.
% Valid values: humanities | demography | methods | hfwg | (unset = default)
%
%   fontset: humanities  -> EB Garamond         (XeLaTeX/LuaLaTeX)
%                           mathpazo + warning  (pdfLaTeX)
%   fontset: demography  -> XITS                (XeLaTeX/LuaLaTeX)
%                           TeX Gyre Termes     (pdfLaTeX)
%   fontset: methods     -> Source Serif 4 + Fira Code (XeLaTeX/LuaLaTeX)
%                           mathpazo + Fira Code or LM Mono (pdfLaTeX)
%   fontset: hfwg        -> XITS, roman headings, CUNY Blue accent (XeLaTeX/LuaLaTeX)
%                           mathpazo            (pdfLaTeX)
%   (default/unset)      -> TeX Gyre Pagella    (XeLaTeX/LuaLaTeX)
%                           mathpazo            (pdfLaTeX)
%
% XeLaTeX/LuaLaTeX fonts are loaded from a local fonts/ directory so the
% bundle is self-contained regardless of system font installation.
% The 'fontpath' YAML variable overrides the default 'fonts/' relative path
% with an absolute path --- required when using the managed template layout
% (e.g., template: ~/Templates/hfwg-tex/hfwgtemplate.tex).
% Example: fontpath: /Users/yourname/Templates/hfwg-tex/fonts/
% =======================================================
\usepackage{iftex}

%% --- FONT PRESET MACRO (set from YAML, used for \ifx branching below) ---
\def\hfwgfontset{hfwg}
\makeatletter
\def\hfwgfontset@h{humanities}
\def\hfwgfontset@d{demography}
\def\hfwgfontset@m{methods}
\def\hfwgfontset@w{hfwg}
%% Boolean flags — set once here, usable everywhere without \makeatletter
\newif\ifhfwghumanities
\newif\ifhfwgdemography
\newif\ifhfwgmethods
\newif\ifhfwghfwg
\ifx\hfwgfontset\hfwgfontset@h\hfwghumanitiestrue\else\hfwghumanitiesfalse\fi
\ifx\hfwgfontset\hfwgfontset@d\hfwgdemographytrue\else\hfwgdemographyfalse\fi
\ifx\hfwgfontset\hfwgfontset@m\hfwgmethodstrue\else\hfwgmethodsfalse\fi
\ifx\hfwgfontset\hfwgfontset@w\hfwghfwgtrue\else\hfwghfwgfalse\fi
\newcommand{\hfwgmethodscodescale}{0.90}

%% --- FONT PATH (override with 'fontpath:' YAML variable for managed templates) ---
\def\hfwgfontsdir{../inst/fonts/}

%% --- FONT LOADING ---

\ifPDFTeX
  %% pdfLaTeX path: uses TeX distribution packages only (no font files needed).
  \usepackage[T1]{fontenc}
  \ifhfwgdemography
    %% demography: TeX Gyre Termes (Times-family); ships with every TeXLive/MiKTeX.
    \usepackage{tgtermes}
  \else
    \usepackage{mathpazo}
    \ifhfwghumanities
      %% humanities: EB Garamond requires XeLaTeX/LuaLaTeX; mathpazo used instead.
      \PackageWarning{hfwgtemplate}{fontset=humanities: EB Garamond requires
        XeLaTeX/LuaLaTeX. Falling back to mathpazo under pdfLaTeX.}
    \fi
    \ifhfwgmethods
      %% methods: try Fira Code for monospace; fall back to Latin Modern Mono.
      \IfFileExists{FiraCode.sty}{%
        \usepackage[scale=\hfwgmethodscodescale]{FiraCode}%
      }{%
        \PackageWarning{hfwgtemplate}{fontset=methods: Fira Code not found in the
          TeX distribution. Using Latin Modern Mono. Install the 'firacode'
          package (TeXLive: tlmgr install firacode).}%
      }
    \fi
    \ifhfwghfwg
      %% hfwg: requires XeLaTeX/LuaLaTeX; mathpazo used instead under pdfLaTeX.
      \PackageWarning{hfwgtemplate}{fontset=hfwg: XITS
        require XeLaTeX/LuaLaTeX. Falling back to mathpazo under pdfLaTeX.}
    \fi
  \fi
\else
  %% XeLaTeX / LuaLaTeX path: uses fontspec with bundled fonts from fonts/ dir.
  %% Each preset uses \IfFileExists to fail gracefully if fonts/ is not populated.
  %% Run scripts/get-fonts.sh to download font files into fonts/.
  \usepackage{fontspec}
  \ifhfwghumanities
    \IfFileExists{\hfwgfontsdir EBGaramond-Regular.otf}{%
      \setmainfont{EBGaramond}[
        Path           = \hfwgfontsdir,
        UprightFont    = *-Regular,
        BoldFont       = *-Bold,
        ItalicFont     = *-Italic,
        BoldItalicFont = *-BoldItalic,
        Extension      = .otf]%
    }{%
      \PackageWarning{hfwgtemplate}{fontset=humanities: \hfwgfontsdir
        EBGaramond-Regular.otf not found. Run scripts/get-fonts.sh, or set
        fontpath: in YAML if using a managed template. Falling back to default.}%
      \IfFontExistsTF{TeX Gyre Pagella}{%
        \setmainfont{TeX Gyre Pagella}%
      }{%
        \PackageWarning{hfwgtemplate}{TeX Gyre Pagella not found; using Latin Modern Roman.}%
        \setmainfont{Latin Modern Roman}%
      }%
    }
  \else\ifhfwgdemography
    \IfFileExists{\hfwgfontsdir XITS-Regular.otf}{%
      \setmainfont{XITS}[
        Path           = \hfwgfontsdir,
        UprightFont    = *-Regular,
        BoldFont       = *-Bold,
        ItalicFont     = *-Italic,
        BoldItalicFont = *-BoldItalic,
        Extension      = .otf]%
    }{%
      \PackageWarning{hfwgtemplate}{fontset=demography: \hfwgfontsdir
        XITS-Regular.otf not found. Run scripts/get-fonts.sh, or set
        fontpath: in YAML if using a managed template. Falling back to default.}%
      \IfFontExistsTF{TeX Gyre Pagella}{%
        \setmainfont{TeX Gyre Pagella}%
      }{%
        \PackageWarning{hfwgtemplate}{TeX Gyre Pagella not found; using Latin Modern Roman.}%
        \setmainfont{Latin Modern Roman}%
      }%
    }
  \else\ifhfwgmethods
    \IfFileExists{\hfwgfontsdir SourceSerif4-Regular.otf}{%
      \setmainfont{SourceSerif4}[
        Path           = \hfwgfontsdir,
        UprightFont    = *-Regular,
        BoldFont       = *-Bold,
        ItalicFont     = *-It,
        BoldItalicFont = *-BoldIt,
        Extension      = .otf]%
    }{%
      \PackageWarning{hfwgtemplate}{fontset=methods: \hfwgfontsdir
        SourceSerif4-Regular.otf not found. Run scripts/get-fonts.sh, or set
        fontpath: in YAML if using a managed template. Falling back to default.}%
      \IfFontExistsTF{TeX Gyre Pagella}{%
        \setmainfont{TeX Gyre Pagella}%
      }{%
        \PackageWarning{hfwgtemplate}{TeX Gyre Pagella not found; using Latin Modern Roman.}%
        \setmainfont{Latin Modern Roman}%
      }%
    }
    \IfFileExists{\hfwgfontsdir FiraCode-Regular.ttf}{%
      \setmonofont{FiraCode}[
        Path        = \hfwgfontsdir,
        UprightFont = *-Regular,
        BoldFont    = *-Bold,
        Scale       = \hfwgmethodscodescale,
        Extension   = .ttf]%
    }{%
      \PackageWarning{hfwgtemplate}{fontset=methods: \hfwgfontsdir
        FiraCode-Regular.ttf not found. Run scripts/get-fonts.sh to install
        Fira Code. Using default monospace font.}%
    }
  \else\ifhfwghfwg
    %% hfwg: XITS body + roman headings for a sharper, classic report style.
    \IfFileExists{\hfwgfontsdir XITS-Regular.otf}{%
      \setmainfont{XITS}[
        Path           = \hfwgfontsdir,
        UprightFont    = *-Regular,
        BoldFont       = *-Bold,
        ItalicFont     = *-Italic,
        BoldItalicFont = *-BoldItalic,
        Extension      = .otf]%
    }{%
      \PackageWarning{hfwgtemplate}{fontset=hfwg: \hfwgfontsdir
        XITS-Regular.otf not found. Run scripts/get-fonts.sh, or set
        fontpath: in YAML if using a managed template. Falling back to default.}%
      \IfFontExistsTF{TeX Gyre Pagella}{%
        \setmainfont{TeX Gyre Pagella}%
      }{%
        \PackageWarning{hfwgtemplate}{TeX Gyre Pagella not found; using Latin Modern Roman.}%
        \setmainfont{Latin Modern Roman}%
      }%
    }
  \else
    %% Default or unrecognized fontset: TeX Gyre Pagella with graceful fallback.
    \ifx\hfwgfontset\empty\else
      \PackageWarning{hfwgtemplate}{Unknown fontset value '\hfwgfontset'.
        Valid values: humanities, demography, methods, hfwg. Using default.}
    \fi
    \IfFontExistsTF{TeX Gyre Pagella}{%
      \setmainfont{TeX Gyre Pagella}%
    }{%
      %% ISSUE 1 FIX: emit warning when TeX Gyre Pagella is unavailable.
      \PackageWarning{hfwgtemplate}{TeX Gyre Pagella not found on this system.
        Using Latin Modern Roman as fallback. Install tex-gyre for the intended
        Palatino-family appearance.}%
      \setmainfont{Latin Modern Roman}%
    }
  \fi\fi\fi\fi
\fi
\makeatother

\usepackage{microtype}
\usepackage{amsmath,amssymb}
\providecommand{\citeproc}[2]{#2}

% =======================================================
% 1b. COLOR PALETTE
% Loaded early so geometry, titlesec, and hyperref can all use hfwgaccent.
% The 'accent' YAML variable overrides all fontset defaults (hex, no '#').
%   accent: E71939   # QC Red
%   accent: 003DA5   # CUNY Blue (hfwg default)
% Fontset defaults (when accent is unset):
%   methods:        deep tech-blue (#1B4F72) with lighter rule (#2E86C1)
%   hfwg:           CUNY Blue (#003DA5, Pantone 286 C)
%   default/others: black
% =======================================================
\usepackage{xcolor}
  \definecolor{hfwgaccent}{HTML}{003DA5}
  \definecolor{hfwgrule}{HTML}{003DA5}

% =======================================================
% 2. PAGE LAYOUT
% humanities:  wider, book-like margins
% methods:     tighter, dense technical report margins
% hfwg:        compact, branded blog-post margins
% default/demography: standard 1-inch
% =======================================================
\ifhfwghumanities
  \usepackage[left=1.3in,right=1.3in,top=1.1in,bottom=1.1in]{geometry}
\else\ifhfwgmethods
  \usepackage[left=0.9in,right=0.9in,top=0.9in,bottom=0.9in]{geometry}
\else\ifhfwghfwg
  \usepackage[left=1.15in,right=1.15in,top=1.1in,bottom=1.1in]{geometry}
\else
  \usepackage[left=1in,right=1in,top=1in,bottom=1in]{geometry}
\fi\fi\fi
\usepackage{setspace}

  %% humanities: slightly more generous leading; methods: tighter; hfwg: balanced
  \ifhfwghumanities
    \setstretch{1.55}
  \else\ifhfwgmethods
    \setstretch{1.25}
  \else\ifhfwghfwg
    \setstretch{1.45}
  \else
    \onehalfspacing
  \fi\fi\fi

%% humanities: wider paragraph indent for book-like feel
\ifhfwghumanities
  \setlength{\parindent}{18pt}
\else
  \setlength{\parindent}{15pt}  %% applies to methods, hfwg, demography, default
\fi
\setlength{\parskip}{0pt}

% =======================================================
% 3. GRAPHICS + FLOATS
% =======================================================
\usepackage{graphicx}
%% eso-pic is only needed for title-page bars.
%% Loading it conditionally keeps default/humanities free of the dependency.
\ifhfwgdemography\usepackage{eso-pic}\else\ifhfwgmethods\usepackage{eso-pic}\else\ifhfwghfwg\usepackage{eso-pic}\fi\fi\fi
\usepackage{float}
\usepackage{multicol}
\usepackage{wrapfig}
\usepackage{longtable,booktabs}
\usepackage{pdflscape}
\usepackage{array}
\usepackage{caption}

\captionsetup{
  font=small,
  labelfont=bf
}

\setcounter{topnumber}{2}
\setcounter{bottomnumber}{1}
\setcounter{totalnumber}{3}

\renewcommand{\arraystretch}{1.1}

% More deterministic float spacing
\setlength{\intextsep}{0.75\baselineskip}
\setlength{\textfloatsep}{1.25\baselineskip}

% --- Auto-scale images ---
\makeatletter
\def\maxwidth{\ifdim\Gin@nat@width>\linewidth\linewidth\else\Gin@nat@width\fi}
\def\maxheight{\ifdim\Gin@nat@height>\textheight\textheight\else\Gin@nat@height\fi}
\makeatother
\setkeys{Gin}{width=\maxwidth,height=\maxheight,keepaspectratio}

% =======================================================
% 4. TABLE AUTOMATION (Pandoc longtable)
% =======================================================
\usepackage{etoolbox}

\AtBeginEnvironment{longtable}{
  \small
  \singlespacing
  \setlength{\LTleft}{\fill}
  \setlength{\LTright}{\fill}
  \setlength{\LTcapwidth}{\textwidth}
  %% Treat special characters as ordinary text inside table cells.
  %% Pandoc does not escape these when converting markdown tables,
  %% so variable names (underscores), exponents (carets), and
  %% percentage values (percent signs) cause LaTeX errors or silent
  %% data loss. Catcode 12 (other) makes them literal characters
  %% without affecting math mode or comments outside this environment.
  \catcode95=12    %% underscore  (_): subscript operator
  \catcode94=12    %% caret      (^): superscript operator
  \catcode37=12    %% percent    (%): comment char -- silent data loss if unescaped
}

\newcommand{\hfwgcompactcolumnssetup}{%
  \setlength{\columnsep}{0.24in}%
  \raggedcolumns
  \footnotesize
  \setstretch{0.98}%
  \setlength{\parindent}{0.85em}%
  \setlength{\parskip}{0pt}%
  \titleformat{\section}{\normalsize\bfseries\color{hfwgaccent}}{}{0pt}{}%
  \titleformat{\subsection}{\footnotesize\bfseries}{}{0pt}{}%
  \titlespacing{\section}{0pt}{0.5em}{0.14em}%
  \titlespacing{\subsection}{0pt}{0.38em}{0.1em}%
}


% =======================================================
% 4b. SPECIAL CHARACTER HANDLING IN BODY TEXT
% =======================================================
%% \usepackage{underscore} makes _ safe in ordinary text mode
%% (e.g. variable names mentioned inline in prose) while
%% preserving its subscript role inside math environments.
%% This complements the longtable catcode fix above.
\usepackage{underscore}

% =======================================================
% 5. SECTION FORMATTING + NUMBERING TOGGLE
%
% default:    bold, no rule
% humanities: small-caps, generous spacing
% demography: bold + thin rule below (journal style)
% methods:    bold + accent color + rule below (futuristic)
% hfwg:       roman bold + CUNY Blue accent + rule below (branded)
% =======================================================
\usepackage{titlesec}

%% Named rule commands for \titleformat after-code.
%% Necessary because \titleformat's optional [after-code] argument uses [],
%% and LaTeX's optional-arg scanner doesn't handle nested [] depth, so
%% \titleformat{...}{}{}{}{}[\titlerule[0.5pt]] would close the outer []
%% at the first ] it encounters. Wrapping in a command avoids nesting.
\newcommand{\hfwgdemorule}{\titlerule[0.3pt]}
\newcommand{\hfwgmetrule}{\color{hfwgrule}\titlerule[0.5pt]}
\newcommand{\hfwgwgrule}{\color{hfwgrule}\titlerule[0.4pt]}

  \setcounter{secnumdepth}{0}
  \setcounter{tocdepth}{0}
  \ifhfwghumanities
    \titleformat{\section}{\large\scshape}{}{0pt}{}
    \titleformat{\subsection}{\normalsize\bfseries\itshape}{}{0pt}{}
    \titleformat{\subsubsection}{\normalsize\itshape}{}{0pt}{}
    \titlespacing{\section}{0pt}{2.5ex plus 0.5ex}{-0.5ex}
    \titlespacing{\subsection}{0pt}{1.5ex plus 0.25ex}{-0.5ex}
    \titlespacing{\subsubsection}{0pt}{1.25ex plus 0.25ex}{-0.5ex}
  \else\ifhfwgdemography
    \titleformat{\section}{\large\bfseries}{}{0pt}{}[\hfwgdemorule\vspace{0.4em}]
    \titleformat{\subsection}{\normalsize\bfseries}{}{0pt}{}
    \titleformat{\subsubsection}{\normalsize\bfseries\itshape}{}{0pt}{}
    \titlespacing{\section}{0pt}{2ex plus 0.5ex}{-0.5ex}
    \titlespacing{\subsection}{0pt}{1.5ex plus 0.25ex}{-0.5ex}
    \titlespacing{\subsubsection}{0pt}{1.25ex plus 0.25ex}{-0.5ex}
  \else\ifhfwgmethods
    \titleformat{\section}{\large\bfseries\color{hfwgaccent}}{}{0pt}{}[\hfwgmetrule\vspace{0.4em}]
    \titleformat{\subsection}{\normalsize\bfseries\color{hfwgaccent}}{}{0pt}{}
    \titleformat{\subsubsection}{\normalsize\itshape\color{hfwgaccent}}{}{0pt}{}
    \titlespacing{\section}{0pt}{1.75ex plus 0.5ex}{-0.5ex}
    \titlespacing{\subsection}{0pt}{1.25ex plus 0.25ex}{-0.5ex}
    \titlespacing{\subsubsection}{0pt}{1ex plus 0.25ex}{-0.5ex}
  \else\ifhfwghfwg
    \titleformat{\section}{\large\bfseries\color{hfwgaccent}}{}{0pt}{}[\hfwgwgrule\vspace{0.4em}]
    \titleformat{\subsection}{\normalsize\bfseries\color{hfwgaccent}}{}{0pt}{}
    \titleformat{\subsubsection}{\normalsize\itshape}{}{0pt}{}
    \titlespacing{\section}{0pt}{2ex plus 0.5ex}{-0.5ex}
    \titlespacing{\subsection}{0pt}{1.5ex plus 0.25ex}{-0.5ex}
    \titlespacing{\subsubsection}{0pt}{1.25ex plus 0.25ex}{-0.5ex}
  \else
    \titleformat{\section}{\large\bfseries}{}{0pt}{}
    \titleformat{\subsection}{\normalsize\bfseries}{}{0pt}{}
    \titleformat{\subsubsection}{\normalsize\bfseries\itshape}{}{0pt}{}
    \titlespacing{\section}{0pt}{2ex plus 0.5ex}{-0.5ex}
    \titlespacing{\subsection}{0pt}{1.5ex plus 0.25ex}{-0.5ex}
    \titlespacing{\subsubsection}{0pt}{1.25ex plus 0.25ex}{-0.5ex}
  \fi\fi\fi\fi

% =======================================================
% 6. HYPERREF + HEADERS/FOOTERS
% =======================================================
%% methods + hfwg: colored links reinforce the accent palette
\ifhfwgmethods
  \usepackage[colorlinks=true,linkcolor=hfwgaccent,urlcolor=hfwgrule,citecolor=hfwgaccent]{hyperref}
\else\ifhfwghfwg
  \usepackage[colorlinks=true,linkcolor=hfwgaccent,urlcolor=hfwgrule,citecolor=hfwgaccent]{hyperref}
\else
  \usepackage[hidelinks]{hyperref}
\fi\fi

%% PDF document metadata — populated from YAML front matter.
%% Enables proper indexing in repositories, readers, and search engines.
%% Multiple authors are semicolon-separated (PDF/XMP convention).
\hypersetup{
  pdftitle    = {Using the HFWG TeX Template from R Markdown: A gentle
introduction for new R users},
  pdfauthor   = {Ada Template},
  pdfsubject  = {A gentle introduction for new R users},
  pdfkeywords = {}
}

\usepackage{fancyhdr}

% =======================================================
% 7. OPTIONAL MODES
% =======================================================
\newif\ifhfwganon
\hfwganonfalse

\newcommand{\hfwgmaincolumns}{1}


% =======================================================
% 8. PANDOC SUPPORT
% =======================================================
\providecommand{\tightlist}{%
  \setlength{\itemsep}{0pt}%
  \setlength{\parskip}{0pt}%
}

\providecommand{\pandocbounded}[1]{#1}

\usepackage{color}
\usepackage{fancyvrb}
\newcommand{\VerbBar}{|}
\newcommand{\VERB}{\Verb[commandchars=\\\{\}]}
\DefineVerbatimEnvironment{Highlighting}{Verbatim}{commandchars=\\\{\}}
% Add ',fontsize=\small' for more characters per line
\usepackage{framed}
\definecolor{shadecolor}{RGB}{248,248,248}
\newenvironment{Shaded}{\begin{snugshade}}{\end{snugshade}}
\newcommand{\AlertTok}[1]{\textcolor[rgb]{0.94,0.16,0.16}{#1}}
\newcommand{\AnnotationTok}[1]{\textcolor[rgb]{0.56,0.35,0.01}{\textbf{\textit{#1}}}}
\newcommand{\AttributeTok}[1]{\textcolor[rgb]{0.13,0.29,0.53}{#1}}
\newcommand{\BaseNTok}[1]{\textcolor[rgb]{0.00,0.00,0.81}{#1}}
\newcommand{\BuiltInTok}[1]{#1}
\newcommand{\CharTok}[1]{\textcolor[rgb]{0.31,0.60,0.02}{#1}}
\newcommand{\CommentTok}[1]{\textcolor[rgb]{0.56,0.35,0.01}{\textit{#1}}}
\newcommand{\CommentVarTok}[1]{\textcolor[rgb]{0.56,0.35,0.01}{\textbf{\textit{#1}}}}
\newcommand{\ConstantTok}[1]{\textcolor[rgb]{0.56,0.35,0.01}{#1}}
\newcommand{\ControlFlowTok}[1]{\textcolor[rgb]{0.13,0.29,0.53}{\textbf{#1}}}
\newcommand{\DataTypeTok}[1]{\textcolor[rgb]{0.13,0.29,0.53}{#1}}
\newcommand{\DecValTok}[1]{\textcolor[rgb]{0.00,0.00,0.81}{#1}}
\newcommand{\DocumentationTok}[1]{\textcolor[rgb]{0.56,0.35,0.01}{\textbf{\textit{#1}}}}
\newcommand{\ErrorTok}[1]{\textcolor[rgb]{0.64,0.00,0.00}{\textbf{#1}}}
\newcommand{\ExtensionTok}[1]{#1}
\newcommand{\FloatTok}[1]{\textcolor[rgb]{0.00,0.00,0.81}{#1}}
\newcommand{\FunctionTok}[1]{\textcolor[rgb]{0.13,0.29,0.53}{\textbf{#1}}}
\newcommand{\ImportTok}[1]{#1}
\newcommand{\InformationTok}[1]{\textcolor[rgb]{0.56,0.35,0.01}{\textbf{\textit{#1}}}}
\newcommand{\KeywordTok}[1]{\textcolor[rgb]{0.13,0.29,0.53}{\textbf{#1}}}
\newcommand{\NormalTok}[1]{#1}
\newcommand{\OperatorTok}[1]{\textcolor[rgb]{0.81,0.36,0.00}{\textbf{#1}}}
\newcommand{\OtherTok}[1]{\textcolor[rgb]{0.56,0.35,0.01}{#1}}
\newcommand{\PreprocessorTok}[1]{\textcolor[rgb]{0.56,0.35,0.01}{\textit{#1}}}
\newcommand{\RegionMarkerTok}[1]{#1}
\newcommand{\SpecialCharTok}[1]{\textcolor[rgb]{0.81,0.36,0.00}{\textbf{#1}}}
\newcommand{\SpecialStringTok}[1]{\textcolor[rgb]{0.31,0.60,0.02}{#1}}
\newcommand{\StringTok}[1]{\textcolor[rgb]{0.31,0.60,0.02}{#1}}
\newcommand{\VariableTok}[1]{\textcolor[rgb]{0.00,0.00,0.00}{#1}}
\newcommand{\VerbatimStringTok}[1]{\textcolor[rgb]{0.31,0.60,0.02}{#1}}
\newcommand{\WarningTok}[1]{\textcolor[rgb]{0.56,0.35,0.01}{\textbf{\textit{#1}}}}

% =======================================================
% 8b. CODE BLOCK STYLING
%
% Applies to:
%   Shaded/Highlighting  — knitr/Pandoc code INPUT blocks
%   verbatim             — knitr console OUTPUT blocks
%
% Design:
%   Input:  blue-tinted background + 3pt left accent rule, footnotesize
%   Output: neutral gray background, footnotesize, no page-split
%
% Implementation notes:
%   - framed is already loaded by highlighting-macros (via snugshade).
%   - \MakeFramed with a custom \FrameCommand draws a colored vrule then
%     a \colorbox; MakeFramed passes the framed content as the argument
%     to the last token of \FrameCommand.
%   - Font size for Highlighting uses \fvset, NOT \RecustomVerbatimEnvironment.
%     fancyvrb v4+ (MiKTeX/TeX Live 2024+) pre-defines \VerbBar internally;
%     \RecustomVerbatimEnvironment re-triggers that initialisation, causing
%     a "Command \VerbBar already defined" conflict with Pandoc's
%     highlighting-macros.  \fvset sets options without re-registering
%     the environment and avoids the conflict entirely.
%   - Font size for verbatim uses \verbatim@font, the LaTeX hook that
%     controls the typeface inside all verbatim environments, for the
%     same reason.
%   - \surroundwithmdframed wraps verbatim externally so its verbatim
%     catcode regime is unaffected, then applies the shaded background.
%   - nobreak=true on the verbatim frame prevents output blocks from
%     splitting across pages; for input blocks, an outer \minipage
%     wrapper achieves the same (blocks longer than \textheight still
%     overflow, which is the acceptable edge-case trade-off).
% =======================================================
\ifhfwgmethods

  %% --- colors ---
  \definecolor{codeinputbg}{HTML}{EBF0F7}   %% blue-tinted input background
  \definecolor{codeoutputbg}{HTML}{F2F2F2}  %% neutral gray output background

  %% --- code INPUT: Shaded + left accent rule, no page split ---
  %% The outer minipage keeps the block together; MakeFramed draws the
  %% colored vrule + background.  Long blocks that exceed \textheight
  %% will still overflow (acceptable trade-off vs. splitting mid-block).
  \renewenvironment{Shaded}{%
    \vspace{3pt}%
    \noindent\begin{minipage}{\linewidth}%
    \def\FrameCommand{%
      \hspace{1pt}%
      {\color{hfwgrule}\vrule width 3pt}%
      \hspace{4pt}%
      \colorbox{codeinputbg}%
    }%
    \MakeFramed{\advance\hsize by -\width \FrameRestore}%
    \footnotesize
    \setlength{\parskip}{0pt}%
    \setlength{\partopsep}{0pt}%
  }{%
    \endMakeFramed
    \end{minipage}%
    \vspace{3pt}%
  }

  %% Fira Code at footnotesize for highlighted code blocks.
  %% \fvset is safe with fancyvrb v4+; \RecustomVerbatimEnvironment is not.
  \fvset{fontsize=\footnotesize}

  %% footnotesize for standard verbatim (console output).
  %% \verbatim@font is the LaTeX hook called inside every verbatim environment
  %% to set the typeface; overriding it here is fancyvrb-independent.
  \makeatletter
  \renewcommand{\verbatim@font}{\footnotesize\ttfamily}
  \makeatother

  %% --- code OUTPUT: shaded verbatim ---
  %% mdframed is intentionally avoided: it loads pgf/tikz, which patches
  %% \@ifnextchar globally and causes infinite recursion ("TeX capacity
  %% exceeded [input stack size=10000]") when hyperref's \href is called
  %% inside \hfwgauthorblock.  The framed package (already loaded above by
  %% highlighting-macros) is sufficient and carries no such risk.
  %%
  %% etoolbox's \AtBeginEnvironment runs before verbatim's catcode changes,
  %% and \AtEndEnvironment runs after they are restored, so \MakeFramed and
  %% \endMakeFramed execute in normal LaTeX mode on either side of the
  %% verbatim content.
  \AtBeginEnvironment{verbatim}{%
    \vspace{3pt}%
    \def\FrameCommand{\fboxsep=4pt\colorbox{codeoutputbg}}%
    \MakeFramed{\FrameRestore}%
  }
  \AtEndEnvironment{verbatim}{%
    \endMakeFramed
    \vspace{3pt}%
  }

\fi

\ifhfwgmethods\else
  %% Non-methods presets use manuscript spacing in prose, but code needs to
  %% remain compact and readable.  Use a smaller monospace face with tight
  %% leading so long lines are less likely to overflow the shaded box.
  \fvset{fontsize=\footnotesize}
  \makeatletter
  \renewcommand{\verbatim@font}{\footnotesize\ttfamily}
  \makeatother
  \AtBeginEnvironment{Shaded}{%
    \footnotesize
    \setstretch{1.05}%
    \setlength{\parskip}{0pt}%
    \setlength{\partopsep}{0pt}%
  }
  \AtBeginEnvironment{verbatim}{%
    \footnotesize
    \setstretch{1.05}%
    \setlength{\parskip}{0pt}%
  }
\fi

%% NOTE (Issue 5): knitr/Quarto may inject booktabs, float, longtable, or
%% other packages already loaded above via header-includes. This is benign ---
%% LaTeX's \usepackage deduplicates package loading --- but may produce harmless
%% "already loaded" notes in verbose log output. No action required.


% --- Pandoc table width arithmetic support ---
\usepackage{calc}
\providecommand{\real}[1]{#1}

% =======================================================
% RUNNING TITLE + PAGESTYLE
% =======================================================
\pagestyle{fancy}
\setlength{\headheight}{14.5pt}
\fancyhf{}
\fancyfoot[C]{\thepage}

% Ensure "plain" pages (e.g., first page after \section in some classes) match
\fancypagestyle{plain}{
  \fancyhf{}
  \fancyfoot[C]{\thepage}
}

%% methods + hfwg: tint the header rule to match the accent color palette.
%% \headrule is the fancyhdr hook that draws the line below the header;
%% redefining it with \color{hfwgrule} keeps the rule weight unchanged
%% while coordinating it with section rules and link colors.
\ifhfwgmethods
  \renewcommand{\headrule}{%
    {\color{hfwgrule}\hrule width\headwidth height\headrulewidth}%
    \vskip-\headrulewidth%
  }
\else\ifhfwghfwg
  \renewcommand{\headrule}{%
    {\color{hfwgrule}\hrule width\headwidth height\headrulewidth}%
    \vskip-\headrulewidth%
  }
\fi\fi

\fancyhead[R]{\small\itshape HFWG: Using the HFWG TeX Template}

% =======================================================
% 9. CSL REFERENCES (Pandoc)
% =======================================================
\newlength{\cslhangindent}
\setlength{\cslhangindent}{1.5em}

\providecommand{\citeproctext}{}

\newenvironment{CSLReferences}[2]
 {%
  %% ISSUE 3 FIX: revert to single column before references when maincolumns: 2.
  %% \twocolumn persists through the body into the bibliography section otherwise.
  %% \onecolumn forces a page break (standard behavior for two-column papers).
  \ifnum\hfwgmaincolumns=2\relax\onecolumn\fi
  \section*{References}
  \begin{list}{}{%
   \setlength{\itemindent}{0pt}
   \setlength{\leftmargin}{0pt}
   \setlength{\parsep}{0pt}
   \ifodd #1
    \setlength{\leftmargin}{\cslhangindent}
    \setlength{\itemindent}{-\cslhangindent}
   \fi
   \ifnum #2 > 0
    \setlength{\itemsep}{#2\baselineskip}
   \else
    \setlength{\itemsep}{0pt}
   \fi
  }%
  \renewcommand{\bibitem}[2][]{\item}
 }
 {\end{list}}

% =======================================================
% 10. R / SOFTWARE MACROS
% =======================================================
\newcommand{\pkg}[1]{\texttt{#1}}
\newcommand{\CRANpkg}[1]{\texttt{#1}}
\newcommand{\code}[1]{\texttt{#1}}

% =======================================================
% BEGIN DOCUMENT
% =======================================================
\begin{document}

% =============================================================
% TITLE PAGE / SNAPSHOT HEADER
% methods:  centered, with thick accent rules top and bottom
% others:   centered (humanities/demography/default)
% =============================================================

%% Shared author-block macro.
%%
%% Author names come from the for(author) loop — Quarto passes author.name
%% through its normalize-author.lua filter. All other author sub-fields
%% (department, institution, email, orcid, etc.) are stripped by that filter
%% and never reach the Pandoc template.
%%
%% Solution: supply contact details via top-level author_* YAML variables,
%% which Quarto does not process and therefore passes through unchanged.
%%
%%   author_dept:   Department of Sociology
%%   author_inst:   Queens College, City University of New York
%%   author_addr:   65-30 Kissena Blvd., Queens, NY 11367, USA
%%   author_email:  hfwg@qc.cuny.edu
%%   author_web:    https://hfwg.qc.cuny.edu
%%   author_orcid:  0000-0002-6197-4453
%%
\newcommand{\hfwgauthorblock}{%
\ifhfwganon\else
%
%% Object author: name + per-author sub-fields (R Markdown / direct Pandoc).
%% NOTE: email/web sub-fields are used WITHOUT an extra \href wrapper.
%% Pandoc processes YAML sub-field values as Markdown and autolinks bare
%% addresses, so wrapping again would cause nested \href recursion.
{\setlength{\parskip}{1pt}%
\textbf{Ada Template}\par
{\small%
%
%
%
%
%
%
}}%
%
\fi
}

\newcommand{\hfwgseriesblock}{%
{\small\setlength{\parskip}{1pt}%
Household Finance Lab\par%
%
}%
}

\ifhfwgmethods
%% -------------------------------------------------------
%% METHODS: centered title page with accent rules
%% Top and bottom bars at 1in from page edge.
%% All content vertically centered between the two bars.
%% -------------------------------------------------------
\newgeometry{top=1in,bottom=1in,left=1.15in,right=1.15in}
\begin{titlepage}
\thispagestyle{empty}
\AddToShipoutPictureBG*{%
  \AtPageUpperLeft{%
    \put(\LenToUnit{\dimexpr 1in+\oddsidemargin\relax},\LenToUnit{-1in}){%
      \raisebox{-4pt}[0pt][0pt]{%
        \textcolor{hfwgaccent}{\rule{\textwidth}{4pt}}%
      }%
    }%
  }%
  \AtPageLowerLeft{%
    \put(\LenToUnit{\dimexpr 1in+\oddsidemargin\relax},\LenToUnit{1in}){%
      \textcolor{hfwgaccent}{\rule{\textwidth}{2pt}}%
    }%
  }%
}%
%% Fixed upper title placement: the top of the title block sits about one
%% top-margin below the accent rule, while the remaining metadata breathes below.
\vspace*{0.65in}%
\begin{center}
\begin{singlespace}
{\setstretch{0.9}\LARGE\bfseries Using the HFWG TeX Template from R
Markdown\par}

\vspace{0.35em}
{\large\color{hfwgaccent} A gentle introduction for new R users\par}
\end{singlespace}
\end{center}
\vfill
\begin{center}
\begin{singlespace}

\hfwgauthorblock

\par\vspace{2\baselineskip}
\hfwgseriesblock

\ifhfwganon\else
{\small May 19, 2026}\par
\fi

\vspace{2\baselineskip}
\textcolor{hfwgrule}{\rule{0.35\textwidth}{0.5pt}}\par
\vspace{0.5em}
\begin{minipage}{0.9\linewidth}
\small\singlespacing
\textbf{Abstract}\par\vspace{0.25em}
This post explains what the HFWG TeX template does and how it fits into
an R Markdown workflow. It is written for students who know the basics
of R and R Markdown but have never worked directly with TeX or LaTeX.
\end{minipage}


\end{singlespace}
\end{center}
\vfill\null
\end{titlepage}
\restoregeometry

\else\ifhfwghfwg
%% -------------------------------------------------------
%% HFWG: one-page blog/post layout.
%% Generic public-facing posts use a designed masthead, compact columns,
%% inline figures, author box, and footer. Use methods for working papers.
%% -------------------------------------------------------
\newgeometry{top=0.55in,bottom=0.55in,left=0.6in,right=0.6in}
\thispagestyle{empty}
\begingroup
\setlength{\parindent}{0pt}
\setlength{\parskip}{0pt}
\singlespacing
\footnotesize
\noindent\textcolor{hfwgaccent}{\rule{\textwidth}{2.2pt}}\par
\vspace{0.45em}

\noindent
\begin{minipage}[t]{0.18\textwidth}
\vspace{0pt}
\fcolorbox{hfwgaccent}{white}{%
  \begin{minipage}[c][0.42in][c]{0.86\linewidth}
  \centering\tiny\bfseries\color{hfwgaccent}HFWG
  \end{minipage}%
}
\end{minipage}\hfill
\begin{minipage}[t]{0.78\textwidth}
\vspace{0pt}
{\small\sffamily\bfseries\color{hfwgaccent}HFWG Blog}\hfill
\ifhfwganon\else
{\small May 19, 2026}%
\fi
\par\vspace{0.35em}

\begingroup
\setbox0=\hbox{\large\bfseries Using the HFWG TeX Template from R
Markdown}%
\ifdim\wd0>\linewidth
  \noindent\resizebox{\linewidth}{!}{\large\bfseries Using the HFWG TeX
Template from R Markdown}\par
\else
  {\large\bfseries Using the HFWG TeX Template from R Markdown\par}%
\fi
\endgroup
\end{minipage}\hfill
\par\vspace{0.35em}

\noindent{\scriptsize\bfseries
Ada Template%
}\par
{\scriptsize Household Finance Lab}\par

\vspace{0.45em}
{\small\bfseries\color{hfwgaccent} A gentle introduction for new R
users\par}

\vspace{0.35em}
\begin{minipage}{0.78\textwidth}
\scriptsize\singlespacing
This post explains what the HFWG TeX template does and how it fits into
an R Markdown workflow. It is written for students who know the basics
of R and R Markdown but have never worked directly with TeX or LaTeX.
\end{minipage}

\vspace{0.55em}
\noindent\textcolor{hfwgrule}{\rule{\textwidth}{0.7pt}}\par
\vspace{0.45em}

\else
%% -------------------------------------------------------
%% DEFAULT / HUMANITIES / DEMOGRAPHY: distinct title pages
%% -------------------------------------------------------
\ifhfwghumanities
\begin{titlepage}
\thispagestyle{empty}
\vspace*{0.08\textheight}
\noindent\rule{0.42\textwidth}{0.4pt}\par
\vspace{2.2em}
{\setstretch{0.92}\huge\scshape Using the HFWG TeX Template from R
Markdown\par}
\vspace{0.35em}
{\Large\itshape A gentle introduction for new R users\par}
\vspace{2.5em}
\begin{singlespace}
\hfwgauthorblock
\par\vspace{1.6\baselineskip}
\hfwgseriesblock
\ifhfwganon\else
\vspace{0.4em}
{\small May 19, 2026}\par
\fi
\end{singlespace}
\vfill
\begin{minipage}{0.72\textwidth}
\small\singlespacing
\noindent\rule{0.25\textwidth}{0.35pt}\par
\textsc{Abstract}\par\vspace{0.3em}
This post explains what the HFWG TeX template does and how it fits into
an R Markdown workflow. It is written for students who know the basics
of R and R Markdown but have never worked directly with TeX or LaTeX.
\end{minipage}
\end{titlepage}
\else\ifhfwgdemography
\newgeometry{top=0.9in,bottom=1in,left=0.85in,right=0.85in}
\begin{titlepage}
\thispagestyle{empty}
\AddToShipoutPictureBG*{%
  \AtPageUpperLeft{%
    \put(0,\LenToUnit{-0.45in}){%
      \textcolor{hfwgaccent}{\rule{\paperwidth}{3pt}}%
    }%
  }%
}%
\begingroup
\setlength{\parindent}{0pt}
\setlength{\parskip}{0pt}
\begin{singlespace}

{\small\bfseries\MakeUppercase{Demographic Research Note}%
%
\hfill
May 19, 2026}\par
\vspace{0.45em}
\noindent\rule{\textwidth}{0.35pt}\par

\vspace{0.11\textheight}
{\setstretch{0.86}\huge\bfseries Using the HFWG TeX Template from R
Markdown\par}

\vspace{0.28em}
{\Large A gentle introduction for new R users\par}

\vspace{1.6em}
\begin{minipage}{0.34\textwidth}
\scriptsize\bfseries\MakeUppercase{Authors}\par
\vspace{0.35em}
\scriptsize\bfseries\MakeUppercase{Affiliation}\par
\vspace{0.35em}
\scriptsize\bfseries\MakeUppercase{Document}\par
\end{minipage}\hfill
\begin{minipage}{0.62\textwidth}
\small
\textbf{Ada Template}\par
\vspace{0.35em}
Household Finance Lab\par
\vspace{0.35em}
Demographic Research Note\par
\end{minipage}

\vspace{1.8em}
\noindent\rule{\textwidth}{0.35pt}\par
\vspace{0.65em}
\begin{minipage}{\textwidth}
\small
\textbf{Abstract.} This post explains what the HFWG TeX template does
and how it fits into an R Markdown workflow. It is written for students
who know the basics of R and R Markdown but have never worked directly
with TeX or LaTeX.
\end{minipage}
\vspace{0.65em}
\par\noindent\rule{\textwidth}{0.35pt}

\end{singlespace}
\endgroup
\vfill
\noindent{\scriptsize\MakeUppercase{Demography fontset specimen}}\hfill{\scriptsize hfwgtex}\par
\end{titlepage}
\restoregeometry
\else
\begin{titlepage}
\thispagestyle{empty}

\vspace*{0.14\textheight}
\begin{center}


{\setstretch{0.9}\LARGE\bfseries Using the HFWG TeX Template from R
Markdown\par}

\vspace{0.25em}
{\large A gentle introduction for new R users\par}

\vspace{2em}

\begin{singlespace}
{\normalsize

\hfwgauthorblock

\par\vspace{2\baselineskip}
\hfwgseriesblock

% --- DATE ---
%% NOTE (Issue 4): To preserve a human-readable date string, use inline R:
%%   date: "`r format(Sys.Date(), '%B %d, %Y')`"
\ifhfwganon\else
\vspace{0.5em}
{\small May 19, 2026}\par
\fi

}
\end{singlespace}


\vspace{2em}

\begin{minipage}{0.8\textwidth}
\small\singlespacing
\textbf{Abstract}\par\vspace{0.25em}
This post explains what the HFWG TeX template does and how it fits into
an R Markdown workflow. It is written for students who know the basics
of R and R Markdown but have never worked directly with TeX or LaTeX.
\end{minipage}

\vfill


\end{center}
\end{titlepage}
\fi\fi
\fi\fi

\ifhfwghfwg
\renewenvironment{longtable}[2][]{%
  \par\smallskip\noindent\footnotesize
  \renewcommand{\caption}[2][]{\textbf{##2}}%
  \begin{tabular}{#2}%
}{%
  \end{tabular}\par\smallskip
}
\renewcommand{\endhead}{}
\renewcommand{\endfirsthead}{}
\renewcommand{\endfoot}{}
\renewcommand{\endlastfoot}{}
\makeatletter
\renewenvironment{figure}[1][]{%
  \par\smallskip\noindent\begin{minipage}{\linewidth}\centering
  \def\@captype{figure}%
}{%
  \end{minipage}\par\smallskip
}
\makeatother
\begin{multicols}{2}
\setlength{\columnsep}{0.24in}
\raggedcolumns
\footnotesize
\setstretch{0.98}
\setlength{\parindent}{0.85em}
\setlength{\parskip}{0pt}
\titleformat{\section}{\normalsize\bfseries\color{hfwgaccent}}{}{0pt}{}
\titleformat{\subsection}{\footnotesize\bfseries}{}{0pt}{}
\titlespacing{\section}{0pt}{0.5em}{0.14em}
\titlespacing{\subsection}{0pt}{0.38em}{0.1em}
\else
\clearpage
\fi
\setcounter{page}{1}

% =======================================================
% MAIN BODY
% =======================================================

% Main text column layout toggle
\ifhfwghfwg\else
\ifnum\hfwgmaincolumns=2\relax
  \twocolumn
\fi
\fi

\ifhfwganon
\begingroup
\small
\noindent\textit{This manuscript is set to anonymous mode. Identifying metadata fields are suppressed on the title page.}\par
\endgroup
\vspace{1em}
\fi

\section{What this package is for}\label{what-this-package-is-for}

If you have just learned R Markdown, you already know the basic idea:
you write text, code, tables, and figures in one \texttt{.Rmd} file,
then click \textbf{Knit} to turn that file into a finished document.

The \texttt{hfwgtex} package helps with the last part of that workflow.
It gives R Markdown a polished academic document design, so you do not
have to learn TeX before you can make a professional-looking PDF.

TeX and LaTeX are the typesetting systems that sit underneath many
academic PDFs. They are powerful, but they can be intimidating if you
are new to them. This template hides most of that complexity. You
control the document with familiar YAML fields at the top of your R
Markdown file.

\section{What the template does}\label{what-the-template-does}

The template controls the appearance of the PDF. It sets the page
margins, fonts, title page, headers, section headings, spacing, tables,
figures, citations, and bibliography formatting.

In a normal R Markdown PDF, you might get a plain-looking document with
default formatting. With this package, the same R Markdown source can
become a Lab report, a working paper, a Snapshot, or a WordPress-ready
blogpost companion.

The important point is that your writing still happens in R Markdown.
You write paragraphs in Markdown, run R code chunks, make figures with
R, and cite sources with citation keys. The template handles the visual
presentation.

\section{The basic workflow}\label{the-basic-workflow}

Most users will follow this pattern:

\begin{enumerate}
\def\labelenumi{\arabic{enumi}.}
\tightlist
\item
  Install the package.
\item
  Run \texttt{hfwgtex::hfwg\_use()} in a project folder.
\item
  Write an R Markdown file.
\item
  Add the HFWG template settings in the YAML header.
\item
  Click \textbf{Knit} in RStudio.
\end{enumerate}

The helper function copies the files your project needs:

\begin{Shaded}
\begin{Highlighting}[]
\NormalTok{hfwgtex}\SpecialCharTok{::}\FunctionTok{hfwg\_use}\NormalTok{()}
\end{Highlighting}
\end{Shaded}

That gives your project a local copy of the template, fonts, and
citation style file. After that, your \texttt{.Rmd} can refer to those
files when it knits.

\section{A minimal YAML header}\label{a-minimal-yaml-header}

The YAML header is the block at the top of an R Markdown document
between two lines of three dashes. It tells R Markdown how to build the
output.

Here is a small example:

\begin{Shaded}
\begin{Highlighting}[]
\PreprocessorTok{{-}{-}{-}}
\FunctionTok{title}\KeywordTok{:}\AttributeTok{ }\StringTok{"My First HFWG Report"}
\FunctionTok{subtitle}\KeywordTok{:}\AttributeTok{ }\StringTok{"A short example"}
\FunctionTok{author}\KeywordTok{:}
\AttributeTok{  }\KeywordTok{{-}}\AttributeTok{ }\FunctionTok{name}\KeywordTok{:}\AttributeTok{ }\StringTok{"Your Name"}

\FunctionTok{institution}\KeywordTok{:}\AttributeTok{ }\StringTok{"Household Finance Lab"}
\FunctionTok{series}\KeywordTok{:}\AttributeTok{ }\StringTok{"Research Note"}
\FunctionTok{number}\KeywordTok{:}\AttributeTok{ }\DecValTok{1}

\FunctionTok{fontset}\KeywordTok{:}\AttributeTok{ hfwg}
\FunctionTok{fontpath}\KeywordTok{:}\AttributeTok{ fonts/}

\FunctionTok{output}\KeywordTok{:}
\AttributeTok{  }\FunctionTok{pdf\_document}\KeywordTok{:}
\AttributeTok{    }\FunctionTok{template}\KeywordTok{:}\AttributeTok{ hfwgtemplate.tex}
\AttributeTok{    }\FunctionTok{latex\_engine}\KeywordTok{:}\AttributeTok{ xelatex}
\PreprocessorTok{{-}{-}{-}}
\end{Highlighting}
\end{Shaded}

You do not need to understand every line at first. The most important
pieces are \texttt{title}, \texttt{author}, \texttt{fontset}, and the
\texttt{output} block.

\section{Choosing a fontset}\label{choosing-a-fontset}

A fontset is a preset design style. It changes the document's typography
and layout.

{\def\LTcaptype{} % do not increment counter
\begin{longtable}[]{@{}ll@{}}
\toprule\noalign{}
Fontset & Best use \\
\midrule\noalign{}
\endhead
\bottomrule\noalign{}
\endlastfoot
\texttt{hfwg} & Household Finance Lab reports and public-facing
documents \\
\texttt{demography} & Journal-style papers using a Times-like look \\
\texttt{humanities} & Essays or theory-heavy writing with a book-like
feel \\
\texttt{methods} & Technical reports with code and compact spacing \\
unset & Plain journal-submission style \\
\end{longtable}
}

For most Lab documents, start with:

\begin{Shaded}
\begin{Highlighting}[]
\FunctionTok{fontset}\KeywordTok{:}\AttributeTok{ hfwg}
\end{Highlighting}
\end{Shaded}

\section{Writing text and running
code}\label{writing-text-and-running-code}

Inside the body of the file, you write ordinary R Markdown. Headings use
\texttt{\#}, emphasis uses \texttt{*italic*} or \texttt{**bold**}, and R
code goes inside code chunks.

For example, this code creates a small figure:

\begin{Shaded}
\begin{Highlighting}[]
\NormalTok{wealth }\OtherTok{\textless{}{-}} \FunctionTok{c}\NormalTok{(}\DecValTok{42000}\NormalTok{, }\DecValTok{91000}\NormalTok{, }\DecValTok{265000}\NormalTok{, }\DecValTok{410000}\NormalTok{)}
\FunctionTok{names}\NormalTok{(wealth) }\OtherTok{\textless{}{-}} \FunctionTok{c}\NormalTok{(}\StringTok{"No diploma"}\NormalTok{, }\StringTok{"High school"}\NormalTok{, }\StringTok{"BA"}\NormalTok{, }\StringTok{"Graduate"}\NormalTok{)}

\FunctionTok{barplot}\NormalTok{(}
\NormalTok{  wealth,}
  \AttributeTok{col =} \FunctionTok{c}\NormalTok{(}\StringTok{"\#B8C7D9"}\NormalTok{, }\StringTok{"\#8FB3D9"}\NormalTok{, }\StringTok{"\#0066CC"}\NormalTok{, }\StringTok{"\#004C99"}\NormalTok{),}
  \AttributeTok{border =} \ConstantTok{NA}\NormalTok{,}
  \AttributeTok{las =} \DecValTok{1}\NormalTok{,}
  \AttributeTok{ylab =} \StringTok{"Median net worth"}\NormalTok{,}
  \AttributeTok{cex.axis =} \FloatTok{0.75}\NormalTok{,}
  \AttributeTok{cex.lab =} \FloatTok{0.85}\NormalTok{,}
  \AttributeTok{cex.names =} \FloatTok{0.8}
\NormalTok{)}
\end{Highlighting}
\end{Shaded}

\begin{figure}
\centering
\pandocbounded{\includegraphics[keepaspectratio,alt={Example bar chart produced from R code.}]{intro-to-hfwgtex-blogpost_files/intro-to-hfwgtex-blogpost_files/figure-latex/example-figure-1.pdf}}
\caption{Example bar chart produced from R code.}
\end{figure}

When you knit the document, R runs the code, creates the figure, and
places it in the PDF.

\section{Blogposts and WordPress}\label{blogposts-and-wordpress}

This package can also help when a Lab document needs to become a
WordPress post on the CUNY Academic Commons.

For that workflow, use \texttt{hfwgtex::blogpost\_pdf}:

\begin{Shaded}
\begin{Highlighting}[]
\FunctionTok{output}\KeywordTok{:}
\AttributeTok{  hfwgtex:}\FunctionTok{:blogpost\_pdf}\KeywordTok{:}
\AttributeTok{    }\FunctionTok{template}\KeywordTok{:}\AttributeTok{ hfwgtemplate.tex}
\AttributeTok{    }\FunctionTok{latex\_engine}\KeywordTok{:}\AttributeTok{ xelatex}
\AttributeTok{    }\FunctionTok{wordpress}\KeywordTok{:}\AttributeTok{ html}
\AttributeTok{    }\FunctionTok{wordpress\_assets}\KeywordTok{:}\AttributeTok{ }\CharTok{true}
\AttributeTok{    }\FunctionTok{wordpress\_checklist}\KeywordTok{:}\AttributeTok{ }\CharTok{true}
\end{Highlighting}
\end{Shaded}

When you click \textbf{Knit}, the package creates several files:

{\def\LTcaptype{} % do not increment counter
\begin{longtable}[]{@{}ll@{}}
\toprule\noalign{}
File & Purpose \\
\midrule\noalign{}
\endhead
\bottomrule\noalign{}
\endlastfoot
\texttt{.pdf} & The polished PDF version \\
\texttt{-wordpress.html} & HTML you can paste into WordPress \\
\texttt{-wordpress-assets/} & Images used by the WordPress companion \\
\texttt{-wordpress-checklist.txt} & A short publishing checklist \\
\end{longtable}
}

The HTML output is designed to be conservative. It does not assume that
special WordPress plugins are active. In most cases, the research
assistant can upload the image files to the Media Library, paste the
HTML into WordPress, preview the post, and make final edits there.

\section{Snapshots}\label{snapshots}

Snapshots are a special kind of Lab publication. They are short,
public-facing, and built around one featured visualization.

For Snapshots, use \texttt{hfwgtex::snapshot\_pdf} and add
\texttt{snapshot:\ true}:

\begin{Shaded}
\begin{Highlighting}[]
\FunctionTok{snapshot}\KeywordTok{:}\AttributeTok{ }\CharTok{true}
\FunctionTok{snapshot\_label}\KeywordTok{:}\AttributeTok{ }\StringTok{"Empirical Snapshot"}
\FunctionTok{snapshot\_feature}\KeywordTok{:}\AttributeTok{ }\StringTok{"figures/my{-}chart.png"}

\FunctionTok{output}\KeywordTok{:}
\AttributeTok{  hfwgtex:}\FunctionTok{:snapshot\_pdf}\KeywordTok{:}
\AttributeTok{    }\FunctionTok{template}\KeywordTok{:}\AttributeTok{ hfwgtemplate.tex}
\AttributeTok{    }\FunctionTok{latex\_engine}\KeywordTok{:}\AttributeTok{ xelatex}
\AttributeTok{    }\FunctionTok{wordpress}\KeywordTok{:}\AttributeTok{ html}
\end{Highlighting}
\end{Shaded}

The Snapshot PDF uses a compact first-page header instead of a full
academic title page. That makes it better suited for short public
reports.

\section{What you should not worry
about}\label{what-you-should-not-worry-about}

You do not need to edit the \texttt{.tex} file directly. You also do not
need to learn LaTeX commands to use the template well.

Most of the time, your job is to:

\begin{enumerate}
\def\labelenumi{\arabic{enumi}.}
\tightlist
\item
  Write clearly in R Markdown.
\item
  Put accurate metadata in the YAML header.
\item
  Make readable figures and tables.
\item
  Knit the document.
\item
  Inspect the PDF and WordPress companion before sharing.
\end{enumerate}

The template is there to make the final document look consistent with
the Lab's style.

\section{A good mental model}\label{a-good-mental-model}

Think of your \texttt{.Rmd} file as the source document. It contains
your words, your code, and your document settings.

Think of the HFWG template as the design system. It decides how your
source document should look when it becomes a PDF.

Think of the WordPress companion as a publishing handoff. It gives you a
web-friendly version of the same material so you do not have to rebuild
the post by hand.

That is the main promise of this package: write once in RStudio, then
produce the formats the Lab needs.


\ifhfwghfwg
\par\vspace{0.65em}
\noindent\textcolor{hfwgrule}{\rule{\linewidth}{0.45pt}}\par
\vspace{0.25em}
{\scriptsize\bfseries\color{hfwgaccent} About the Author}\par
\vspace{0.2em}
\begin{minipage}{\linewidth}
\scriptsize\singlespacing
\begin{minipage}[c]{0.20\linewidth}
\colorbox{hfwgaccent!10}{\rule{0pt}{0.58in}\rule{0.58in}{0pt}}
\end{minipage}\hfill
\begin{minipage}[c]{0.74\linewidth}
\textbf{Ada Template}\par
Research Associate, Household Finance Lab\par
Household Finance Lab\par

Creates practical guides and short public posts for Household Finance
Lab readers.
\end{minipage}
\end{minipage}\par
\end{multicols}
\vfill
\noindent\textcolor{hfwgrule}{\rule{\textwidth}{0.7pt}}\par
\vspace{0.25em}
\noindent
\begin{minipage}[t]{0.24\textwidth}
\scriptsize\bfseries\color{hfwgaccent}HFWG Blog
\end{minipage}\hfill
\begin{minipage}[t]{0.72\textwidth}
\scriptsize\raggedleft
Household Finance Lab. May 19, 2026.
\end{minipage}\par
\endgroup
\restoregeometry
\fi

% =======================================================
% REFERENCES
% =======================================================

\end{document}
